\documentclass[troisieme]{fiche}
\metadonnees{7}{Les loups de Seeonee}{Bloc B — Ce qu'on peut, ce qu'on doit}

\begin{document}
\entete

\objectifs{%
\textbf{Langue} : \emph{can}, \emph{could}, \emph{be able to} ; la demande et la
permission.\par
\textbf{Texte} : Rudyard Kipling, \og Mowgli's Brothers \fg, \emph{The Jungle Book}
(1894).\par
\textbf{Durée} : 1 h — exercices 20 min, texte et questions 25 min, version 15 min.}

%% ===================================================================
\exercice[12 min]{\emph{Can}, \emph{could}, \emph{be able to}}

\rappel{\textbf{Rappel.} \emph{Can} dit ce qu'on sait ou ce qu'on peut faire maintenant,
\emph{could} au passé. Là où \emph{can} n'a pas de forme — au futur, après \emph{to} —, on
emploie \emph{be able to}.}

\consigne{Complétez.}

\begin{anglais}
\begin{enumerate}[label=\arabic*.,itemsep=2.5mm,leftmargin=*]
  \item Father Wolf \trou{can} hear the jackal at the entrance of the cave.
  \item When he was a cub, he \trou{could} not hunt alone.
  \item Tomorrow the cubs \trou{will be able to} leave the cave.
  \item Even the tiger \trou{cannot} catch Tabaqui when he goes mad.
  \item The wolves would like to \trou{be able to} hunt in peace.
  \item Mowgli \trou{could} not speak, but he \trou{could} laugh.
  \item \trou{Can} you see the moon from the mouth of the cave?
  \item Nobody has ever \trou{been able to} tame a jackal.
  \item The little boy \trou{was able to} climb into the cave by himself.
    \emph{(une seule fois, et il a réussi)}
  \item She \trou{could} smell the tiger long before she saw him.
\end{enumerate}
\end{anglais}

\note[La règle la plus utile]{Un modal n'a ni infinitif ni futur : on ne dit pas \emph{to can}
ni \emph{will can}. On passe alors par \emph{be able to} (3, 5, 8). Phrase 9 : pour une
réussite unique dans le passé, l'anglais préfère \emph{was able to} à \emph{could}.}

%% ===================================================================
\exercice[8 min]{Demander, permettre, refuser}

\rappel{\textbf{Rappel.} \emph{Can I\ldots ?} demande la permission, \emph{Can you\ldots ?}
demande un service. \emph{Could} est plus poli que \emph{can}.}

\consigne{Traduisez en anglais.}

\begin{anglais}
\begin{enumerate}[label=\arabic*.,itemsep=3mm,leftmargin=*]
  \item Est-ce que je peux entrer ?
    \lignes{1}
    \corr{Can I come in? \emph{ou} May I come in?}
  \item Pourriez-vous m'aider, s'il vous plaît ?
    \lignes{1}
    \corr{Could you help me, please?}
  \item Tu ne peux pas rester ici.
    \lignes{1}
    \corr{You can't stay here.}
  \item Il ne pouvait pas dormir.
    \lignes{1}
    \corr{He couldn't sleep.}
  \item Sais-tu nager ?
    \lignes{1}
    \corr{Can you swim?}
  \item Nous ne pourrons pas venir demain.
    \lignes{1}
    \corr{We won't be able to come tomorrow.}
\end{enumerate}
\end{anglais}

\note[Un faux ami de construction]{\og Savoir \fg{} au sens de \og être capable de \fg{} se
dit \emph{can} : \emph{can you swim?} et non \emph{do you know to swim?} Et \emph{can} n'est
jamais suivi de \emph{to} : \emph{I can swim}.}

%% ===================================================================
\newpage
\section{Texte}

\noindent\small\textit{Dans les collines de Seeonee, en Inde, une famille de loups s'éveille
au crépuscule.}\normalsize

\begin{texte}
It was seven o'clock of a very warm evening in the Seeonee hills when Father Wolf woke up from
his day's rest, scratched himself, \gl[to yawn]{yawned}{bâiller}, and spread out his
\gl[a paw]{paws}{une patte} one after the other to get rid of the sleepy feeling in their
tips. Mother Wolf lay with her big gray nose dropped across her four tumbling,
\gl[to squeal]{squealing}{couiner} \gl[a cub]{cubs}{un petit (de loup)}, and the moon shone
into the mouth of the cave where they all lived. ``Augrh!'' said Father Wolf. ``It is time to
hunt again.'' He was going to spring down hill when a little shadow with a
\gl[bushy]{bushy}{touffu} tail crossed the \gl[a threshold]{threshold}{un seuil} and
\gl[to whine]{whined}{geindre}: ``Good luck go with you, O Chief of the Wolves. And good luck
and strong white teeth go with noble children that they may never forget the hungry in this
world.''

It was the \gl[a jackal]{jackal}{un chacal} --- Tabaqui, the Dish-licker --- and the wolves of
India \gl[to despise]{despise}{mépriser} Tabaqui because he runs about making
\gl[mischief]{mischief}{des bêtises, du désordre}, and telling tales, and eating
\gl[a rag]{rags}{un chiffon} and pieces of leather from the village rubbish-heaps. But they
are afraid of him too, because Tabaqui, more than anyone else in the jungle, is
\gl[to be apt to]{apt}{avoir tendance à} to go mad, and then he forgets that he was ever
afraid of anyone, and runs through the forest biting everything in his way. Even the tiger
runs and hides when little Tabaqui goes mad, for madness is the most
\gl[disgraceful]{disgraceful}{honteux} thing that can \gl[to overtake]{overtake}{frapper,
s'emparer de} a wild creature.

``Enter, then, and look,'' said Father Wolf \gl[stiffly]{stiffly}{avec raideur}, ``but there
is no food here.''

``For a wolf, no,'' said Tabaqui, ``but for so \gl[mean]{mean}{humble, misérable} a person as
myself a dry bone is a good feast. Who are we, the jackal people, to pick and choose?'' He
\gl[to scuttle]{scuttled}{détaler} to the back of the cave, where he found the bone of a
\gl[a buck]{buck}{un chevreuil} with some meat on it, and sat cracking the end merrily.

``All thanks for this good meal,'' he said, licking his lips. ``How beautiful are the noble
children! How large are their eyes! And so young too!''

Now, Tabaqui knew as well as anyone else that there is nothing so unlucky as to compliment
children to their faces. It pleased him to see Mother and Father Wolf look uncomfortable.
\end{texte}
\source{Rudyard Kipling, \og Mowgli's Brothers \fg, \emph{The Jungle Book}, Londres,
Macmillan, 1894.}

\partiel[15 min]{Questions}
\consigne{Répondez aux deux premières questions en français, aux deux dernières en anglais.}

\begin{enumerate}[label=\arabic*.,itemsep=2.5mm,leftmargin=*]
  \item Que fait la famille de loups au début du texte ?
    \lignes{2}
    \corr{Le père loup se réveille de son repos de la journée, se gratte, bâille et étire ses
    pattes ; la mère est couchée, le museau posé en travers de ses quatre petits qui se
    bousculent et couinent. Le père s'apprête à partir chasser.}
  \item Pourquoi les loups méprisent-ils Tabaqui, et pourquoi le craignent-ils quand même ?
    \lignes{3}
    \corr{Ils le méprisent parce qu'il fait des bêtises, rapporte des histoires et mange des
    chiffons et des bouts de cuir dans les tas d'ordures des villages. Ils le craignent parce
    qu'il est plus sujet que tout autre à devenir fou, et qu'il mord alors tout ce qu'il
    rencontre.}
  \item \en{Tabaqui says the children are beautiful. Why is that not a compliment?}
    \lignes{3}
    \corr{Because he knows perfectly well that praising children to their faces brings bad
    luck. He says it on purpose, and he enjoys seeing the parents uncomfortable. The
    politeness is the attack.}
  \item \en{The story is told in the past, but one paragraph is in the present. Find it and
    say why.}
    \lignes{2}
    \corr{Le deuxième paragraphe : \emph{the wolves of India despise Tabaqui}, \emph{he runs
    about}, \emph{even the tiger runs and hides}. Ce n'est plus le récit d'un soir mais une
    vérité générale sur les chacals, vraie aujourd'hui encore.}
\end{enumerate}

\note[Une politesse qui n'en est pas]{\emph{Good luck go with you, O Chief of the Wolves} :
Tabaqui parle comme à la cour. Tout le passage joue sur l'écart entre les mots, très polis,
et ce que chacun sait de celui qui les prononce.}

%% ===================================================================
\section{Version}
\consigne{Traduisez le deuxième paragraphe du texte, de \emph{It was the jackal} à
\emph{a wild creature}.}

\lignes{12}

\corr{\textbf{Proposition de traduction.} C'était le chacal — Tabaqui, le Lécheur
d'écuelles — et les loups de l'Inde méprisent Tabaqui parce qu'il court partout à faire des
bêtises, à colporter des histoires, et à manger des chiffons et des morceaux de cuir dans les
tas d'ordures des villages. Mais ils le craignent aussi, car Tabaqui, plus que quiconque dans
la jungle, est sujet à devenir fou ; il oublie alors qu'il a jamais eu peur de quoi que ce
soit, et court à travers la forêt en mordant tout ce qui se trouve sur son passage. Même le
tigre s'enfuit et se cache quand le petit Tabaqui devient fou, car la folie est la chose la
plus honteuse qui puisse frapper une créature sauvage.}

\note[Choix de traduction 1]{\emph{because he runs about making mischief, and telling tales,
and eating rags} : trois gérondifs coordonnés après \emph{runs about}. Le français les rend
par des infinitifs introduits par \og à \fg, ou par \og en \fg{} + participe.}

\note[Choix de traduction 2]{\emph{is apt to go mad} : \emph{to be apt to} = \og avoir
tendance à \fg, \og être sujet à \fg. Rien à voir avec \og apte \fg, qui se dit \emph{fit}
ou \emph{able}.}

\note[Choix de traduction 3]{\emph{the most disgraceful thing that can overtake a wild
creature} : après un superlatif, le français emploie le subjonctif — \og la chose la plus
honteuse qui \emph{puisse} frapper \fg.}

%% ===================================================================
\partiel{Lexique à mémoriser}
\begin{multicols}{2}\small
\begin{itemize}[label=--,itemsep=0.8mm,leftmargin=*]
  \item \textbf{a cave} — une grotte
  \item \textbf{a cub} — un petit (de loup, d'ours)
  \item \textbf{a paw} — une patte
  \item \textbf{to yawn} — bâiller
  \item \textbf{to hunt} — chasser
  \item \textbf{to despise} — mépriser
  \item \textbf{mischief} — les bêtises, le désordre
  \item \textbf{leather} — le cuir
  \item \textbf{mad} — fou ; \textit{et} furieux
  \item \textbf{to hide} — se cacher ; hid, hidden
  \item \textbf{a bone} — un os
  \item \textbf{a meal} — un repas
\end{itemize}
\end{multicols}

\end{document}
